\section{Introduction}
\label{sec:intro}

A results section can contain hundreds of numbers: rates, intervals, differences, counts and
p-values. In most projects each one is copied by hand from a script's output into the text. The
copy is correct on the day it is made. It stops being correct when the analysis is rerun with a
fixed bug, a changed exclusion or new data, and nothing in the manuscript records which output a
number came from, so finding every affected sentence depends on the author's memory. Text edited
by an author or by a writing assistant adds a second route for error, since a number can be
rounded differently, attached to the wrong arm of a comparison, or carried over from an earlier
draft.

Two families of tools address part of this. Literate-programming systems such as
Sweave~\cite{leisch2002sweave} and its successors run the analysis inside the document, so numbers
cannot drift, but they require the analysis to be written in the document's language and rerun at
every build. Many projects keep analysis and manuscript apart: the analysis is a set of scripts
that write tables, and the paper is written in LaTeX or a word processor. Checking tools take the
finished text as given. statcheck~\cite{nuijten2016prevalence} extracts test statistics reported
in APA style and recomputes their p-values. sciwrite-lint~\cite{samsonau2026sciwrite} checks
citations and uses a language model to compare numbers in the text with the paper's own typeset
tables. Neither looks at the analysis output, so a table that is itself out of date passes.

\pv{} works from the tables the analysis writes. It turns each table cell into a LaTeX macro and
the manuscript refers to cells by address, so rerunning the analysis and rebuilding the paper
updates every number. A linter reports numbers still typed by hand. The harder problem is an
existing manuscript, where someone has to decide which cell each typed number reports. That is a
reading task: ``the 2~\AA{} padding raised Vina's top-1 success by 6.8 points'' identifies an
engine, an arm, a metric and a statistic, and the table has many cells containing 6.8. We give
that reading to a language model, but only as a proposal. The program looks up the proposed cell
and accepts it if, and only if, its value formatted as the text formats it reproduces the written
number. A wrong proposal costs coverage and cannot make a wrong number look right.

This paper describes the design (Section~\ref{sec:design}), the evaluation on one manuscript
(Sections~\ref{sec:evaluation} and~\ref{sec:results}), and what that evaluation does and does not
show (Section~\ref{sec:discussion}). The manuscript is the DockPrepBench
benchmark~\cite{alobo2026dockprepbench}, a study of how ligand and receptor preparation choices
change docking success on the PoseBusters set~\cite{buttenschoen2024posebusters}. Its results
section is dense with paired differences, intervals and corrected p-values, and the tables behind
it are generated by scripts, which is the setting \pv{} is built for.
